\chapter*{General Conclusion}
\addcontentsline{toc}{chapter}{General Conclusion}

This internship set out to design and build a real-time anomaly detection platform for branch teller operations at Amen Bank, covering \emph{retrait}, \emph{versement}, \emph{virement}, and \emph{remise de ch\`eques}. The result is a working end-to-end system that ingests transaction events, computes behavioral features, scores them with two complementary deep learning models, applies regulatory rules, and surfaces explainable alerts to branch supervisors through an interactive dashboard.

The platform monitors two independent actors --- the client initiating the transaction and the employee processing it --- reflecting the dual-actor nature of teller operations where anomalies can originate from either side. A streaming architecture on RedPanda chains four microservices through message topics, supported by a hot/cold data layer and a batch pipeline for nightly profile reconstruction. The Autoencoder detects point anomalies while the LSTM identifies sequential deviations, with a sigmoid-weighted fusion that adapts to account maturity. Every alert carries a human-readable SHAP explanation from deterministic templates, letting supervisors act on the system's output without machine learning expertise.

Beyond the detection pipeline, the project invested heavily in the infrastructure that makes a machine learning system maintainable. MLflow tracks every experiment. An automated promotion gate ensures only models matching or exceeding current production performance are deployed. Three independent retraining triggers --- a weekly schedule, a Prometheus-driven drift detector, and a supervisor feedback divergence monitor --- close the loop between serving and improvement. Structured JSON logging, Prometheus metrics, and Grafana dashboards provide the observability needed to diagnose issues without inspecting individual events.

The system was validated on a synthetic dataset of 243{,}000 events generated from five archetypes calibrated against Amen Bank's published data. The validation confirms that the pipeline works end-to-end and that the retraining subsystem can detect drift and update models autonomously. The synthetic data caveat is explicit: these results demonstrate engineering readiness, not production-grade detection accuracy. The architecture is designed so that replacing synthetic data with real transactions requires no code changes --- the same feature computation, the same models, and the same retraining pipeline apply without modification.

Two lessons stand out. Technically, the most impactful decision in the project was not a model choice but a feature engineering one: replacing profile dumps with contrast features raised AUC from 0.57 to over 0.98 with no change to architecture or hyperparameters, and enforcing a single pure feature function across streaming, batch, and simulation eliminated an entire class of train--serve skew. More broadly, building a model turned out to be a small fraction of the effort --- the majority of the engineering work lies in the infrastructure that keeps that model correct, observable, and improvable over time.
